\section{Discussion}
\label{sec:discussion}

The eight experiments together support three substantive claims about
how real markets generate the price patterns that empirical research
has found contradictory.

\subsection{Time-on-market is the missing third dimension}
\label{sec:disc-tom}

Runs~3 and~4 show that two distinct manipulations --- raising seller
patience and lowering agent density --- both produce stable
deviations from the textbook equilibrium price. Their effect signs
are opposite (Run~3 overshoots, Run~4 undershoots), but their
mechanism is the same: both alter the typical \emph{time on market}
of an unsold home. Patience increases the upper bound on time on
market directly; density increases it indirectly by decreasing
encounter rates per unit time. In each case the mismatch between the
balancer's price-to-quantity correction and the underlying microstructural
adjustment shows up as a persistent equilibrium offset.

This suggests that the textbook supply-and-demand pair $\hat{Q}_{s}(p)$
and $\hat{Q}_{d}(p)$ omits a state variable. A more accurate
characterisation would be $\hat{Q}_{s}(p, \tau)$ and
$\hat{Q}_{d}(p, \tau)$ where $\tau$ is expected time on market.
Empirically, real-estate datasets that include days-on-market
information vastly outperform those that do not for short-run
forecasting; our model offers a generative explanation for that
finding.

\subsection{Bid-selection asymmetry, not estimation bias, drives bubbles}
\label{sec:disc-bias}

The estimation function \cref{eq:estimation} has zero mean error: an
agent's expected valuation is exactly the home's current value. Yet
the EDEM speculative regime (Runs 6--8) produces unbounded upward
drift even with this unbiased estimator. The mechanism is not
behavioural bias but a statistical asymmetry: each seller selects the
\emph{maximum} bid received, and the maximum of $N$ noisy estimates
of a fair value is, in expectation, larger than the fair value
itself. \cref{eq:value-update} compounds this asymmetry
multiplicatively across epochs.

This refines the standard ``divergence of opinion'' story.
\citet{miller1977}'s static premium emerges from constrained supply:
the $N$ shares are sold to the $N$ most optimistic of $K > N$
investors. EDEM's dynamic premium emerges from \emph{repeated}
constrained supply: each epoch's winners feed back into the next
epoch's anchor value. The static premium is a corner of the dynamic
premium, recovered when the feedback loop is severed.

A practical implication is that empirical premium-vs-discount
findings depend strongly on the sample's position along this
feedback trajectory, which short-window studies cannot identify.
Two studies that find opposite signs may both be correct
characterisations of distinct epochs in the same underlying process.

\subsection{Mean-reverting balancers cannot fully cancel the asymmetry}
\label{sec:disc-balancer}

Run~7's three-way comparison shows that no setting of the EDEM
balancer coefficient $C_{b}$ within the parameter grid we examined
restores $v_{t}/v^{*} = 1$. The mean-reverting case ($C_{b} = +1$)
plateaus near $4 \times$; the trend-following case ($C_{b} = -1$)
continues drifting upward; the no-balancer case dominates both. A
crucial corollary is that \emph{policy interventions cast as
balancer adjustments will not, by themselves, prevent persistent
mispricing.} The asymmetry is structural to the bidding mechanism;
to neutralise it, an intervention would have to change either the
selection rule (e.g., uniform random matching instead of max-bid) or
the per-epoch update rule (e.g., a winsorised mean instead of the
arithmetic mean). \cref{sec:extensions} sketches both alternatives.

The companion observation --- that the no-balancer case dominates
the trend-following case in our simulation --- refines the framing
in \citet{arbuzov2018edem}, which described the trend-following
``Silicon Valley real estate'' regime as the most explosive. That
framing turns out to overstate the role of the balancer's sign
relative to its presence. We attribute the discrepancy to the
agent-pool depletion that strong $|C_{b}|$ induces: trend-following
removes sellers, but the remaining sellers are all that the
multiplicative update rule can act on.

\subsection{Implications for valuation algorithms}
\label{sec:disc-ml}

The same mechanism that drives EDEM's exponential bubble in Run~6
operates in any pricing system whose inputs are anchored to recent
market-clearing prices. A machine-learning real-estate valuation
algorithm trained on historical sale prices estimates not the home's
fair value but the typical winning bid for a similar home --- which
is, by construction, the maximum of the relevant noisy estimates,
and therefore systematically above the fair value. Deploying such an
algorithm at scale risks providing the market with a coordinating
signal that pulls realised winning bids upward, which then feeds
back into the next training cycle.

We do not advance EDEM as an explanation for any specific real-world
mispricing event. The 2021 wind-down of Zillow's algorithmic-iBuying
programme \citep{zillow2021ibuying} is one episode in which a
historical-sale-price-anchored valuation system at scale incurred
substantial losses; the public record is consistent with EDEM's
order-statistic prediction but is also consistent with several other
explanations (timing of housing-market regime changes, operational
mispricing, capital-allocation constraints), and we have not
attempted to discriminate between them. The robust theoretical
finding is more modest: an unbiased ML valuation trained on
historical clearing prices is, by construction, fitting the
upper-order statistic rather than the population mean of plausible
sale prices, and at deployment time it functions as a
positively-biased estimator. A defensible deployment would correct
for this --- for example by predicting the median rather than the
mean of plausible sale prices, or by using a censored-data
likelihood that accounts for the offer-acceptance threshold ---
rather than treating the historical record as a clean training
target.

\subsection{Limitations}
\label{sec:disc-limits}

Several modelling choices restrict the scope of the findings.

\begin{itemize}[leftmargin=*]
\item \textbf{Single asset class.} EDEM models a single class of
      indistinguishable homes; portfolio-level effects (cross-asset
      hedging, sector rotation) are absent.
\item \textbf{No leverage.} Buyers carry no debt and sellers no
      mortgages; financialised real-estate markets that exhibited the
      most dramatic anomalies of the past two decades require leverage
      to model accurately \citep{geanakoplos2012leverage,baptista2016boe}.
\item \textbf{No macroeconomic environment.} The supply and demand
      schedules in \cref{eq:Qs-update,eq:Qd-update} respond only to
      price; they do not respond to interest-rate shocks, employment,
      or migration.
\item \textbf{Adaptive vs.\ static estimators.} Agents do not learn
      from sale outcomes. A learning extension --- Bayesian belief
      updating, or reinforcement-learning bidding --- would either
      converge toward bias-corrected estimation (closing the
      bubble channel) or amplify it (opening a stronger one),
      depending on how counterfactual observations are weighted.
\item \textbf{Run~5 timing.} The 12{,}000-tick budget for the single-
      shock scenario is too short for the market to reach its
      post-shock floor before the patience boost fires. Doc~2's
      specific claim that patience returns the price to the new
      textbook equilibrium would require a longer simulation.
      The figure as published makes the more general slow-adjustment
      claim that \cref{sec:bg-dynamics} actually requires.
\end{itemize}
